% numbers.tex: generated by export_numbers.py from results/*.json. Do not edit by hand.
\newcommand{\PH}[1]{\textcolor{red!70!black}{[#1]}}

\newcommand{\nVE}{232}  % analysis.V_EARTH
\newcommand{\nVesc}{544}  % analysis.V_ESC
\newcommand{\nVmax}{776}  % v_esc + v_E
\newcommand{\nVEmin}{217}  % annual minimum of the Earth speed
\newcommand{\nVEmax}{247}  % annual maximum of the Earth speed
\newcommand{\nUMeV}{16.36}  % end of the Bahcall interpolant; verification/PROOF.txt
\newcommand{\nEend}{4.39}  % 2U^2/m_N for 131Xe (keV)
\newcommand{\nEb}{4.639}  % 2U^2/m_N for 124Xe (keV)
\newcommand{\nEendMin}{4.23}  % 2U^2/m_N for 136Xe (keV)
\newcommand{\nMlo}{6.7}  % mass_for(E_end_131, v_max)
\newcommand{\nMloMin}{6.65}  % min over isotopes
\newcommand{\nMloMax}{6.69}  % max over isotopes
\newcommand{\nEc}{0.39}  % E_end (v_E/v_max)^2, 131Xe
\newcommand{\nEcRound}{0.4}  % rounded E_c
\newcommand{\nMhi}{7.7}  % mass_for(0.51 keV, v_E), 131Xe
\newcommand{\nMhiTwo}{7.66}  % same, two decimals
\newcommand{\nMhiThirty}{5.8}  % mass_for(0.30 keV, v_E), 131Xe
\newcommand{\nMhiLight}{7.46}  % 124Xe
\newcommand{\nMhiHeavy}{7.79}  % 136Xe
\newcommand{\nMhiHeavyOne}{7.8}  % 136Xe, one decimal
\newcommand{\nMhiHeavyThirty}{5.893}  % 136Xe at 0.30 keV
\newcommand{\nMhiList}{4.7,\,5.8,\,6.7,\,7.7,\,9.1,\,11.0}  % m_hi at 0.2,0.3,0.4,0.51,0.7,1.0 keV
\newcommand{\nMhiVEmin}{8.2}  % m_hi(0.51) at v_E = 217
\newcommand{\nMhiVEa}{8.2}  % m_hi(0.51) at v_E = 218
\newcommand{\nMhiVEb}{7.1}  % m_hi(0.51) at v_E = 248
\newcommand{\nMhiBinShift}{8.05}  % m_hi(0.56), 131Xe
\newcommand{\nMhiBinShiftHeavy}{8.19}  % m_hi(0.56), 136Xe
\newcommand{\nQstarLo}{3.181534858300}  % verification/certificate.json
\newcommand{\nQstarHi}{3.181534858301}  % verification/certificate.json
\newcommand{\nQstarShort}{3.1815}  % verification/certificate.json, truncated
\newcommand{\nEstarLo}{0.165530074482}  % verification/certificate.json
\newcommand{\nEstarHi}{0.165530074483}  % verification/certificate.json
\newcommand{\nEstarShort}{0.165}  % verification/certificate.json, truncated
\newcommand{\nEstarMax}{0.182}  % verification/PROOF.txt, all isotopes and +-3 sigma columns
\newcommand{\nModeMeV}{6.44}  % mode of the Bahcall spectrum
\newcommand{\nNatMhi}{7.793}  % verification/natural_xenon: m_hi(136Xe)
\newcommand{\nNatMhiLong}{7.7932}  % same, four decimals
\newcommand{\nNatStrictLow}{6.687}  % mass for 124Xe to reach E_b
\newcommand{\nNatMhiThirty}{5.893}  % m_hi(136Xe, 0.30 keV)
\newcommand{\nNatSupportMin}{5.76}  % support bound at 3.5 keV, 124Xe
\newcommand{\nNatBtwoMass}{6.0233}  % B2 counterexample mass
\newcommand{\nNatCells}{39}  % mass cells in the five-term certificate
\newcommand{\nNatCellLo}{6.65}  % partition start in verification/natural_xenon/certify.py (covers all masses reaching E_b)
\newcommand{\nNatCellHi}{7.80}  % partition end, above m_hi(136Xe)
\newcommand{\nNatP}{0.146}  % weighted operator p at 0.51 keV
\newcommand{\nNatR}{0.092}  % weighted operator r at 0.51 keV
\newcommand{\nNatPthirty}{0.113}  % p at 0.30 keV
\newcommand{\nNatRthirty}{0.089}  % r at 0.30 keV
\newcommand{\nNatMarginB}{6.9\times10^{-4}}  % min_i (2 l_i' + p)
\newcommand{\nNatMarginC}{1.2\times10^{-3}}  % -max_i (l_i''+l_i'^2+p l_i'+r)
\newcommand{\nNatLnu}{1.0\times10^{6}}  % L applied to the target at 0.51 keV
\newcommand{\nFFSkin}{0.9}  % Helm skin s (fm) in verification/natural_xenon/certify.py; Lewin & Smith
\newcommand{\nFFSkinStep}{0.1}  % difference between certificate skin 0.9 fm and Engel skin 1 fm
\newcommand{\nFFEngelSkin}{1}  % Engel skin (fm) quoted by Lewin & Smith; form_factor_variations.json
\newcommand{\nFFSkinLow}{0.8}  % certificate skin minus 0.1 fm; form_factor_variations.json helm_skin_0.8
\newcommand{\nFFEngelRms}{0.93}  % Engel r_rms/A^{1/3} (fm); form_factor_variations.json engel
\newcommand{\nFFEngelC}{1.25\times10^{-3}}  % max of l''+l'^2+p l'+r, Engel Helm, 0.51 keV
\newcommand{\nFFEngelCthirty}{5.96\times10^{-4}}  % same maximum at 0.30 keV
\newcommand{\nFFLowqB}{-2.72\times10^{-3}}  % min of 2 l'+p, F^2=1-q^2<r^2>/3, 0.51 keV
\newcommand{\nFFDelta}{1.0\times10^{-6}}  % delta in rechoose(), verification/natural_xenon/certify.py
\newcommand{\nFFEngelP}{0.14197983}  % re-chosen p, Engel, 0.51 keV; form_factor_variations.json
\newcommand{\nFFEngelR}{0.09050014}  % re-chosen r, Engel, 0.51 keV
\newcommand{\nFFEngelMB}{1.01\times10^{-6}}  % min(2 l'+p), Engel re-chosen, 0.51 keV neighborhood
\newcommand{\nFFEngelMC}{1.01\times10^{-6}}  % -max(c), Engel re-chosen, 0.51 keV neighborhood
\newcommand{\nFFEngelLnu}{1.0\times10^{6}}  % L nu, Engel re-chosen, 0.51 keV neighborhood
\newcommand{\nFFEngelPthirty}{0.10877233}  % re-chosen p, Engel, 0.30 keV
\newcommand{\nFFEngelRthirty}{0.08820597}  % re-chosen r, Engel, 0.30 keV
\newcommand{\nFFEngelMBthirty}{1.00\times10^{-6}}  % min(2 l'+p), Engel re-chosen, 0.30 keV
\newcommand{\nFFEngelMCthirty}{1.00\times10^{-6}}  % -max(c), Engel re-chosen, 0.30 keV
\newcommand{\nFFEngelLnuThirty}{6.6\times10^{5}}  % L nu, Engel re-chosen, 0.30 keV
\newcommand{\nFFLowqP}{0.14872490}  % re-chosen p, low-q, 0.51 keV
\newcommand{\nFFLowqR}{0.09883439}  % re-chosen r, low-q, 0.51 keV
\newcommand{\nFFLowqMB}{1.00\times10^{-6}}  % min(2 l'+p), low-q re-chosen, 0.51 keV
\newcommand{\nFFLowqMC}{1.00\times10^{-6}}  % -max(c), low-q re-chosen, 0.51 keV
\newcommand{\nFFLowqLnu}{1.0\times10^{6}}  % L nu, low-q re-chosen, 0.51 keV
\newcommand{\nStructGlo}{4.782}  % floor of the smallest closed-form gap G at 0.51 keV; verification/structural
\newcommand{\nStructGloThirty}{1.963}  % floor of the smallest closed-form gap G at 0.30 keV; verification/structural
\newcommand{\nMomCap}{69}  % sufficient mean-square-radius cap at 0.51 keV (fm^2); verification/structural
\newcommand{\nMomCapThirty}{44}  % sufficient mean-square-radius cap at 0.30 keV (fm^2); verification/structural
\newcommand{\nSkinMax}{0.3}  % neutron-skin hypothesis upper end (fm); verification/structural
\newcommand{\nSkinDq}{0.01}  % local momentum width of the shape-independent skin test (MeV)
\newcommand{\nSkinLoss}{0.255}  % ceil of the fractional curvature loss at 0.51 keV; prose uses below this
\newcommand{\nSkinLossThirty}{0.434}  % ceil of the fractional curvature loss at 0.30 keV
\newcommand{\nSkinRtwo}{27.49}  % ceil of the largest absolute weak mean-square radius (fm^2)
\newcommand{\nEcFullLo}{4.45}  % full-endpoint critical-threshold bracket, lower end (keV)
\newcommand{\nEcFullHi}{4.58}  % full-endpoint critical-threshold bracket, upper end (keV)
\newcommand{\nEcWitness}{7}  % witness mass for the full-endpoint member at nEcFullHi (GeV)
\newcommand{\nEcIsoHi}{0.415}  % ceil of the nine single-isotope critical thresholds (keV)
\newcommand{\nNatLowA}{5.764}  % low-mass functional interval, lower edge (GeV)
\newcommand{\nNatLowB}{5.893}  % low-mass functional interval, upper edge (GeV)
\newcommand{\nSinW}{0.2387}  % sin^2 theta_W in analysis.cevns_spectrum; Q_W = N-(1-4 sin^2 theta_W)Z
\newcommand{\nCresstThr}{0.0301}  % CaWO4 threshold (keV), analysis.CRESST_THRESHOLD
\newcommand{\nCresstEmax}{16}  % CaWO4 ROI upper edge (keV), analysis.CRESST_EMAX
\newcommand{\nCresstBins}{24}  % complement.py geomspace(threshold, emax, 25) has 24 bins
\newcommand{\nCresstRes}{4.6}  % CRESST-III baseline resolution (eV), analysis.CRESST_SIGMA_EV; not smeared again
\newcommand{\nNatEb}{4.639}  % 124Xe end point (keV)
\newcommand{\nNatEendMin}{4.23}  % 136Xe end point (keV)
\newcommand{\nIdealDE}{0.05}  % ideal bin width (keV)
\newcommand{\nIdealNmass}{120}  % ideal mass grid
\newcommand{\nIdealMmin}{4}  % 
\newcommand{\nIdealMmax}{16}  % 
\newcommand{\nIdealLo}{6.68}  % numerical lower edge, 0.51 keV, 1e-12
\newcommand{\nIdealHi}{8.14}  % numerical upper edge, 0.51 keV, 1e-12
\newcommand{\nIdealLoTol}{5.23}  % numerical lower edge, 0.51 keV, 1e-3
\newcommand{\nIdealHiTol}{8.24}  % numerical upper edge, 0.51 keV, 1e-3
\newcommand{\nIdealEmptyUpTo}{0.30}  % largest threshold with empty 1e-12 window
\newcommand{\nIdealOpenFrom}{0.35}  % smallest threshold with non-empty 1e-12 window
\newcommand{\nIdealRabove}{5.5\times10^{-3}}  % residual 1 GeV above m_hi
\newcommand{\nRelLoS}{4.6}  % released window, class S, lower
\newcommand{\nRelLoI}{4.96}  % released window, class I, lower
\newcommand{\nRelLoIone}{5.0}  % released window, class I, lower, one decimal
\newcommand{\nRelHiI}{12.6}  % released window, class I, upper
\newcommand{\nBsupport}{2.8}  % energy above which 1e-3 of 8B events lie (keV)
\newcommand{\nMloPractical}{5.2}  % mass_for(2.8 keV, v_max)
\newcommand{\nEcPractical}{0.25}  % 2.8 keV (v_E/v_max)^2
\newcommand{\nLiveA}{1.17}  % SR0 livetime (t y)
\newcommand{\nLiveB}{2.34}  % SR1 livetime (t y)
\newcommand{\nLive}{3.5}  % SR0+SR1 livetime (t y)
\newcommand{\nFluxPrior}{3.81}  % 8B flux prior (percent)
\newcommand{\nFluxPriorRound}{3.8}  % 8B flux prior, rounded
\newcommand{\nFluxPriorDelta}{0.0381}  % 8B flux prior, fraction
\newcommand{\nBflux}{5.25\times10^{6}}  % SNO 8B flux cm^-2 s^-1, analysis.SOLAR
\newcommand{\nBrate}{3.4}  % 8B per t y in [0.51,5] keV through the response
\newcommand{\nShmDev}{3}  % 6 GeV SHM rate vs release (percent)
\newcommand{\nOwnDev}{0.2}  % own vs released 8B spectrum, 0.5-3 keV (percent)
\newcommand{\nRobVEa}{13.3}  % robustness.json
\newcommand{\nRobVEb}{11.7}  % robustness.json
\newcommand{\nRobMhiVEa}{8.2}  % robustness.json
\newcommand{\nRobMhiVEb}{7.1}  % robustness.json
\newcommand{\nSepMillionBelow}{12.8}  % isotropic separation exposure exceeds 1e6 t y below this mass
\newcommand{\nSepTenMillionBelow}{12.3}  % exceeds 1e7 t y below this mass
\newcommand{\nSepThirteen}{4.7\times10^{5}}  % separation exposure at 13 GeV
\newcommand{\nSepFifteen}{1083}  % separation exposure at 15 GeV
\newcommand{\nSepSixteen}{138}  % separation exposure at 16 GeV
\newcommand{\nSepEighteen}{16.9}  % separation exposure at 18 GeV
\newcommand{\nSepTwenty}{6.9}  % separation exposure at 20 GeV
\newcommand{\nSepHighLo}{2.3}  % separation exposure, 25-40 GeV, min
\newcommand{\nSepHighHi}{2.7}  % separation exposure, 25-40 GeV, max
\newcommand{\nSepMassNow}{23}  % mass reached at the SR0+SR1 exposure
\newcommand{\nSepMassTwenty}{18}  % mass reached at 20 t y
\newcommand{\nSepMassThousand}{15}  % mass reached at 1000 t y
\newcommand{\nSepSabove}{4.7}  % class S: infinite above this mass
\newcommand{\nDiscEpsLow}{3}  % lowest quoted exposure (t y)
\newcommand{\nDiscEpsMax}{3\times10^{4}}  % largest exposure in the scan
\newcommand{\nDiscRatioMinLow}{1.4}  % smallest HI/SHM ratio above the window at 3 t y
\newcommand{\nDiscRatioMaxLow}{666}  % HI/SHM ratio at 8 GeV, 3 t y
\newcommand{\nDiscRatioMidLo}{9}  % 12-14 GeV at 3 t y
\newcommand{\nDiscRatioMidHi}{20}  % 12-14 GeV at 3 t y
\newcommand{\nDiscRatioMidLoHigh}{42}  % 12-14 GeV at 1e4 t y
\newcommand{\nDiscRatioMidHiHigh}{63}  % 12-14 GeV at 1e4 t y
\newcommand{\nDiscNfourteenK}{7}  % fog index, 14 GeV, 1e3 t y
\newcommand{\nDiscNfourteenTenK}{35}  % fog index, 14 GeV, 1e4 t y
\newcommand{\nDiscNsixteenPeak}{7}  % fog index peak, 16 GeV
\newcommand{\nDiscNsixteenPeakEps}{949}  % exposure of that peak (t y)
\newcommand{\nDiscEpsTenK}{9487}  % exposure grid point near 1e4 t y
\newcommand{\nDiscEpsK}{949}  % exposure grid point near 1e3 t y
\newcommand{\nDiscNsixteenTenK}{4}  % fog index, 16 GeV, 1e4 t y
\newcommand{\nDiscSatLo}{1.5\times10^{-46}}  % saturation sigma, 12-14 GeV
\newcommand{\nDiscSatHi}{2.6\times10^{-46}}  % saturation sigma, 12-14 GeV
\newcommand{\nDiscRatioTwentyLo}{1.6}  % 20 GeV ratio range
\newcommand{\nDiscRatioTwentyHi}{2.2}  % 
\newcommand{\nDiscRatioHighLo}{1.4}  % 30-50 GeV ratio range
\newcommand{\nDiscRatioHighHi}{1.7}  % 
\newcommand{\nDiscRatioEightLo}{578}  % 8 GeV ratio range over exposures
\newcommand{\nDiscRatioEightHi}{764}  % 
\newcommand{\nDiscUgrid}{3}  % first positive shell speed (km/s), discovery.py geomspace(3, v_esc, 90)
\newcommand{\nDiscUEight}{3.37}  % dominant shell speed at 8 GeV, 3 t y; discovery.json fstar
\newcommand{\nDiscWEight}{99}  % minimum dominant-shell weight at 8 GeV, percent, over the exposure grid
\newcommand{\nDiscRateEight}{1.28\times10^{-3}}  % 8 GeV f* rate over SHM rate at 3 t y
\newcommand{\nDiscTvEight}{41}  % 8 GeV shape total variation vs 8B, percent of counts, at 3 t y
\newcommand{\nDiscUTwelve}{3.18}  % dominant shell speed at 12 GeV, 1e4 t y
\newcommand{\nDiscWTwelve}{95}  % dominant weight at 12 GeV, 1e4 t y, percent
\newcommand{\nDiscUTwelveFast}{184}  % weight-averaged speed of the remaining 12 GeV shells at 1e4 t y
\newcommand{\nDiscTvTwelve}{4.1}  % 12 GeV f* shape total variation vs 8B, percent, at 1e4 t y
\newcommand{\nDiscUFourteen}{3.00}  % dominant shell speed at 14 GeV, 1e4 t y
\newcommand{\nDiscWFourteen}{98.7}  % dominant weight at 14 GeV, 1e4 t y, percent
\newcommand{\nDiscTvFourteen}{1.0}  % 14 GeV f* shape total variation vs 8B, percent, at 1e4 t y
\newcommand{\nDiscTvThirty}{50}  % 30 GeV escape-shell shape total variation vs 8B, percent of counts
\newcommand{\nCawoRate}{300}  % 8B per t y in CaWO4 (CRESST-III efficiencies)
\newcommand{\nCompSlo}{2.6}  % CaWO4 exposure, class S, min over window masses
\newcommand{\nCompShi}{4.5}  % 
\newcommand{\nCompIlo}{0.07}  % CaWO4 exposure, class I
\newcommand{\nCompIhi}{0.29}  % 
\newcommand{\nCompFactorLo}{12}  % ratio S/I
\newcommand{\nCompFactorHi}{68}  % 
\newcommand{\nBkgRate}{7.5}  % AC+ER+radiogenic NR per combined t y in the ROI; XENONnT release templates (XENON:2024hup), livetime-weighted SR0+SR1
\newcommand{\nBkgCounts}{26.4}  % same backgrounds, expected events in the SR0+SR1 exposure
\newcommand{\nBkgFactorLo}{2.3}  % min over 13-25 GeV of (exposure with bkg)/(8B-only exposure), isotropic class
\newcommand{\nBkgFactorHi}{2.9}  % max of the same ratio
\newcommand{\nBkgSepTwenty}{18.0}  % separation exposure at 20 GeV with fixed backgrounds (t y)
\newcommand{\nBkgSepSixteen}{360}  % separation exposure at 16 GeV with fixed backgrounds (t y)
\newcommand{\nEthrLow}{0.30}  % second threshold of the natural-xenon and structural certificates (keV)
\newcommand{\nEcFullCells}{1355}  % mass cells in the full-endpoint exclusion; verification/structural
\newcommand{\nEcFullMassLo}{6.65}  % full-endpoint exclusion, lower mass (GeV); verification/structural
\newcommand{\nEcFullMassHi}{27}  % full-endpoint exclusion, upper mass (GeV); above it Eq. (closed) at 4.45 keV applies
\newcommand{\nModAmpB}{3.3}  % 8B annual modulation, percent half-range, Kepler orbit e = 0.0167
\newcommand{\nModPhaseB}{3 January}  % date of the 8B maximum (perihelion)
\newcommand{\nModAmpDM}{11}  % smallest modulation over exact isotropic mimics at 7 GeV, ideal 0.51 keV detector, percent
\newcommand{\nModAmpDMNNLS}{13}  % same for the NNLS vertex mimic
\newcommand{\nModAmpDMRel}{20}  % smallest modulation over mimics at 7 GeV, released XENONnT observable, percent
\newcommand{\nModPhaseDM}{2 June}  % date of the mimic's maximum (Earth speed maximum)
\newcommand{\nModRateIdeal}{294}  % 8B events per t y above 0.51 keV with unit efficiency (ideal detector)
\newcommand{\nModExposure}{3.0}  % deviance 9 from timing alone at 7 GeV, ideal detector, worst-case mimic (t y)
\newcommand{\nModExposureLo}{3.3}  % same at 6.7 GeV
\newcommand{\nModExposureHi}{2.5}  % same at 7.6 GeV
\newcommand{\nModExposureRel}{100}  % deviance 9 from timing alone at 7 GeV, released observable, worst-case mimic (t y)
\newcommand{\nModExposureRelLo}{110}  % same at 6.7 GeV
\newcommand{\nModExposureRelHi}{83}  % same at 7.6 GeV
\newcommand{\nProfSepTwenty}{19.4}  % profiled separation exposure at 20 GeV (t y); results/profiled_backgrounds.json
\newcommand{\nProfSepSixteen}{545}  % profiled separation exposure at 16 GeV (t y)
\newcommand{\nProfFactorLo}{1.04}  % min profiled/fixed over 13--25 GeV
\newcommand{\nProfFactorHi}{1.51}  % max profiled/fixed over 13--25 GeV
\newcommand{\nFreeSepTwenty}{24.1}  % free-background-cone separation exposure at 20 GeV (t y)
\newcommand{\nFreeFactorHi}{2.42}  % max free/fixed over 13--25 GeV
\newcommand{\nProfSigAc}{9.04/5.76}  % AC normalization widths for SR0/SR1 (percent); release YAML
\newcommand{\nProfSigEr}{100}  % ER normalization width (percent); release YAML
\newcommand{\nProfSigRg}{58}  % radiogenic normalization width (percent); release YAML
\newcommand{\nProfThetaTwenty}{0.58}  % largest absolute fitted normalization shift at 20 GeV, in prior sigma
\newcommand{\nProfTarget}{9}  % target Asimov deviance
\newcommand{\nProfDevianceFactor}{2}  % Poisson deviance prefactor
\newcommand{\nProfYieldSpan}{1}  % radiogenic yield range in release-node units
\newcommand{\nProfTabThirteenMass}{13}  % profiled-background table, 13 GeV mass
\newcommand{\nProfTabThirteenFixed}{1.2\times10^{6}}  % profiled-background table, 13 GeV fixed
\newcommand{\nProfTabThirteenProfiled}{1.7\times10^{6}}  % profiled-background table, 13 GeV profiled
\newcommand{\nProfTabThirteenFree}{1.7\times10^{6}}  % profiled-background table, 13 GeV free
\newcommand{\nProfTabSixteenMass}{16}  % profiled-background table, 16 GeV mass
\newcommand{\nProfTabSixteenFixed}{360}  % profiled-background table, 16 GeV fixed
\newcommand{\nProfTabSixteenProfiled}{545}  % profiled-background table, 16 GeV profiled
\newcommand{\nProfTabSixteenFree}{613}  % profiled-background table, 16 GeV free
\newcommand{\nProfTabTwentyMass}{20}  % profiled-background table, 20 GeV mass
\newcommand{\nProfTabTwentyFixed}{18.0}  % profiled-background table, 20 GeV fixed
\newcommand{\nProfTabTwentyProfiled}{19.4}  % profiled-background table, 20 GeV profiled
\newcommand{\nProfTabTwentyFree}{24.1}  % profiled-background table, 20 GeV free
\newcommand{\nProfTabTwentyFiveMass}{25}  % profiled-background table, 25 GeV mass
\newcommand{\nProfTabTwentyFiveFixed}{8.0}  % profiled-background table, 25 GeV fixed
\newcommand{\nProfTabTwentyFiveProfiled}{8.3}  % profiled-background table, 25 GeV profiled
\newcommand{\nProfTabTwentyFiveFree}{19.3}  % profiled-background table, 25 GeV free
\newcommand{\nFreeFactorTwentyFive}{2.42}  % free/fixed separation exposure at 25 GeV; results/profiled_backgrounds.json
\newcommand{\nProfPctTwentyFive}{3.6}  % profiled/fixed - 1 at 25 GeV (percent); results/profiled_backgrounds.json
\newcommand{\nCertMassGeV}{122.3}  % nuclear mass used by the convexity certificate (GeV); verification/certificate.json mass_MeV
\newcommand{\nQstarMinus}{3.12436}  % q_* for the -3 sigma column (MeV); verification/certificate.json
\newcommand{\nQstarPlus}{3.23871}  % q_* for the +3 sigma column (MeV); verification/certificate.json
\newcommand{\nUMinus}{16.16}  % end point U of the -3 sigma column (MeV); verification/certificate.json
\newcommand{\nUPlus}{16.56}  % end point U of the +3 sigma column (MeV); verification/certificate.json
\newcommand{\nLogRemExp}{55}  % log-series remainder bound 10^-N; verification/certificate.json
\newcommand{\nQstarWidthExp}{12}  % width of the certified q_* enclosure is 10^-N MeV; verification/certificate.json
\newcommand{\nFiveTermRatio}{0.98}  % sqrt(E0/E_b) of the five-term argument; verification/structural/results.json
\newcommand{\nFiveTermPct}{2}  % 1 - sqrt(E0/E_b), percent of the momentum range
\newcommand{\nTailPeakExp}{11}  % natural-xenon 8B spectrum on [E0, E_b] is below 10^-N of its peak (max ratio 1.10e-12); b8_tail_over_peak
\newcommand{\nShellU}{400}  % Galactic shell speed drawn in Fig. 1; viz/render_shell.py U
\newcommand{\nShellKinkLo}{168}  % |u - v_E| for the Fig. 1 shell
\newcommand{\nShellKinkHi}{632}  % u + v_E for the Fig. 1 shell
\newcommand{\nPubBrateSRzero}{3.6}  % fixed input: published SR0 8B rate (t^-1 y^-1), XENON:2024hup
\newcommand{\nPubBrateSRone}{3.3}  % fixed input: published SR1 8B rate (t^-1 y^-1), XENON:2024hup
